% ECONOMICS_PROBLEM: {"id": "O17-612", "title": "Informality and firm growth", "jel_code": "O17", "source_id": "OP-0034"}
% !TeX program = xelatex
\documentclass[11pt,a4paper]{article}
\usepackage[margin=23mm]{geometry}
\usepackage{fontspec}
\setmainfont{Latin Modern Roman}
\usepackage{amsmath,amssymb,mathtools}
\usepackage{xcolor,enumitem,booktabs,longtable,array,xurl,hyperref}
\definecolor{muted}{HTML}{53636C}
\hypersetup{colorlinks=true,urlcolor=blue,linkcolor=blue}
\setlength{\parindent}{0pt}
\setlength{\parskip}{5pt}
\newcommand{\R}{\mathbb R}
\newcommand{\E}{\mathbb E}
\newcommand{\N}{\mathbb N}
\newcommand{\ind}{\mathbf 1}
\newcommand{\Var}{\operatorname{Var}}
\newcommand{\Cov}{\operatorname{Cov}}
\newcommand{\supp}{\operatorname{supp}}
\DeclareMathOperator*{\argmin}{arg\,min}
\DeclareMathOperator*{\argmax}{arg\,max}
\newcommand{\entrymeta}[3]{{\sffamily\small\color{muted}\textbf{\texttt{#1}}\quad|\quad #2\par #3\par}}
\newcommand{\Problemref}[1]{\texttt{#1}}
\newcommand{\JELref}[2]{\texttt{#2} (JEL \texttt{#1})}
\begin{document}
\section*{Informality and firm growth}
\textbf{Problem ID:} \texttt{O17-612}\quad \textbf{JEL:} \texttt{O17}\par
% BEGIN ECONOMICS STATEMENT
\label{problem:OP-0034}
\label{audited:ECON-071}

\label{atlas:prob:D4}
\noindent\textbf{Original atlas ID:} D4 (entry 34). \textbf{Classification:} Editorial registration-versus-policy-bundle identification problem.

\noindent\textbf{Original atlas question:} Does formalization cause firms to become more productive, or do productive firms simply formalize?

\subsection*{Definitions and assumptions}
Fix baseline enterprises, workers, and horizon $H$. Distinguish active-firm registration $F_{jt}\in\{0,1\}$ from compliant-job status $G_{ijt}\in\{0,1\}$; a registered firm can employ unregistered workers. Define $F_{jt}=0$ after closure, and $G_{ijt}=0$ whenever worker $i$ has no job at firm $j$, including after closure or during nonemployment. Otherwise $G_{ijt}=1$ precisely when that job satisfies the specified compliance rules. A policy vector $\pi$ specifies registration costs, tax schedules, labor requirements, enforcement, and access to services. Let $V_{jt}(\pi)$ be real value added, $L_{jt}(\pi)$ total labor hours, and $S_{jt}(\pi)$ enterprise survival, tracking baseline firms after exit and entry changes. Set output and hours to zero after genuine closure, and report survival separately. A structural model $M$ specifies demand, production, selection into formality, evasion, and market adjustment.
\subsection*{Formal research question}
For announced policy contrasts $\pi_1,\pi_0$, identify
\[
 \Theta_H=\left(
 \mathbb E[V_{jt}(\pi_1)-V_{jt}(\pi_0)],\;
 \mathbb E[S_{jt}(\pi_1)-S_{jt}(\pi_0)],\;
 \Delta F_t,\Delta G_t,\Delta L_t\right)_{t=0}^H,
\]
where $\Delta F_t,\Delta G_t,\Delta L_t$ use specified baseline firm and worker populations. Separately, in a model defining an intervention that changes registration while holding the other policy components fixed, identify or bound the registration effect on investment and productivity. State when that intervention is feasible and when the data only identify a bundled reform. For aggregate productivity, use a defined output-hours or production-function measure with positive denominators and explicit price adjustment.
\subsection*{Interpretation and scope}
Productive firms can select into registration; the observed registered-unregistered gap therefore does not identify formalization's causal effect. A registration-assistance instrument requires exclusion from productivity other than through registration, which can fail if assistance also supplies training or information. Compliance costs, access to finance, and enforcement changes are distinct policy channels. Conditioning on surviving firms can conceal induced exit or entry, and revenue productivity can change merely because prices change. The formal sector's output share, the fraction of registered firms, and the fraction of compliant jobs need not move together. Employment composition and workers' outside options belong to the equilibrium response and cannot be inferred from registry counts alone.
\subsection*{Source audit and status}
[S1] reviews informal firms and [S2] explicitly distinguishes informality margins. [S3], verified as NBER Working Paper 34445, was issued in 2025 and revised in July 2026. It studies minimum wages and informality; it is adjacent evidence, not a direct formalization experiment. The registration-versus-bundled-policy identification problem is editorial, and this audit does not certify universal unresolved status.

\subsection*{References}
\begin{enumerate}[label=\textnormal{[S\arabic*]},leftmargin=*]
\item Rafael La Porta and Andrei Shleifer (2014). \emph{Informality and Development}. Journal of Economic Perspectives 28(3), 109--126. \url{https://doi.org/10.1257/jep.28.3.109}. \textit{Locator:} Primary publisher, working-paper, or author record: bibliographic information and abstract; full-text theorem verification is not claimed. \textit{Role:} Evidence and conceptual distinctions among informal firms.
\item Gabriel Ulyssea (2018). \emph{Firms, Informality, and Development: Theory and Evidence from Brazil}. American Economic Review 108(8). \url{https://doi.org/10.1257/aer.20141745}. \textit{Locator:} Primary publisher, working-paper, or author record: bibliographic information and abstract; full-text theorem verification is not claimed. \textit{Role:} Model distinguishing firm and labor informality; supports separate extensive/intensive margins.
\item Ellora Derenoncourt, François Gerard, Lorenzo Lagos, and Claire Montialoux (2025). \emph{Minimum Wages and Informality}. NBER Working Paper 34445, issued November 2025; revised July 2026. \url{https://www.nber.org/papers/w34445}. \textit{Locator:} Primary publisher, working-paper, or author record: bibliographic information and abstract; full-text theorem verification is not claimed. \textit{Role:} Minimum-wage policy and informal-sector wage spillovers; adjacent evidence rather than direct causal evidence on formalization. Primary record verifies source identity and 2026 revision. The mismatch between wage-policy evidence and formalization effects is flagged in the entry.
\end{enumerate}

\subsection*{Integrated refinement: ECON-071}
{\sffamily\small\color{muted}Informal firms as stepping stones to formal growth\par}
This refinement adopts the additional source conventions
(page~\pageref{audited:conventions}), subject to its local restrictions.
\textbf{Initially informal firms' productive transition paths.}
Track an initially informal cohort through policy-induced registration, exit
and growth. Let $F_{i,t}(z)$ denote status, with an explicit post-closure rule,
and $\Pi_{i,t}(z)$ net profits including exit outcomes. Identify
\[
P_z(h)=\Pr(F_{i,h}(z)=1\mid F_{i,0}=0),\qquad
V_z=\mathbb E\sum_{t=0}^{T}\delta^t\Pi_{i,t}(z).
\]
Declare registration costs, taxes, enforcement and formal-sector benefits as
components of the full regime. Use pretreatment cohort definitions and
consistent real-cost/profit accounting. Distinguish genuinely productive
transitions from relabeling, survival selection and displacement of other firms.
Registration, compliant jobs, sustained productivity and household welfare are
separate outcomes, so administrative counts alone do not establish growth.
\subsection*{Supplied source audit for ECON-071}
Local [S1], [S2], etc. in this audit and the references immediately below
belong to ECON-071, independently of the original entry's reference numbering.
[S1] and [S2] directly support informal firms as potential stepping stones. The fixed-cohort transition/value formulation is editorial; the author's name is Gabriel Ulyssea, despite a spelling inconsistency on a review citation line.
\subsection*{References for integrated refinement ECON-071}
{\footnotesize\raggedright
\par\noindent\textbf{[S1]} Gabriel Ulyssea, Matteo Bobba, Lucie Gadenne, and Mariaflavia Harari (2025). \emph{Informality}. VoxDevLit 6(2), April 2025.\par \textit{Verified locator / source role:} Primary publisher page checked for review identity, authors, version and background. The specific published agenda is corroborated in [S2]; the exact formal target is editorial.\par \url{https://voxdev.org/voxdevlit/informality}\par\smallskip
\par\noindent\textbf{[S2]} Oliver Hanney / VoxDev (living compilation). \emph{Evidence gaps: Key unanswered questions in development economics}. Accessed 8 October 2026. The page currently covers 20 topics; the ``16-topics URL is a historical slug.\par \textit{Verified locator / source role:} Topic heading: Informality. Supports the broad research agenda; this is not a verbatim quotation or an attribution of the displayed mathematical target.\par \url{https://voxdev.org/topic/evidence-gaps-key-unanswered-questions-across-16-topics-development-economics}\par\smallskip
}

\subsection*{Common conventions: Source mathematical conventions}

Unless an entry states otherwise, agent and firm sets are finite. State and action spaces are measurable, strategies and outcome maps are measurable, probability laws are specified, and expectations exist for the targets under discussion. Infinite-horizon discount factors lie in $(0,1)$ and discounted objectives are finite. Norms, tolerances and complexity input sizes must be specified for computational comparisons. Constants or primitives that appear locally are inputs, not empirical facts asserted by this catalogue.

\textbf{Identification.} A statistical model $m\in\mathcal M$ implies an observable law $P_m$ and target $\theta(m)$. At law $P$, its identified set is
\[
\Theta(P;\mathcal M)=\{\theta(m):m\in\mathcal M,\ P_m=P\}.
\]
Point identification holds when this set is a singleton, or equivalently when $P_{m_1}=P_{m_2}$ implies $\theta(m_1)=\theta(m_2)$ for all compatible models. Sharp bounds mean bounds attained, or approached when the boundary is not attained, by compatible models. Writing this set is a definition, not a solution: the substantive task is to establish informative restrictions and derive or estimate the set. An unrestricted-confounding model may give only trivial bounds.

\textbf{Inference.} Identification, estimability and valid uncertainty quantification are separate questions. Fix $\alpha\in(0,1)$. A confidence procedure $C_n$ over a declared law class $\mathcal P$ has asymptotic uniform coverage at level $1-\alpha$ when
\[
\liminf_{n\to\infty}\inf_{P\in\mathcal P}\Pr_P\{\theta(P)\in C_n\}\geq1-\alpha.
\]
For set-identified targets the coverage event must be defined separately, for example coverage of the full identified set. If an entry uses identification-indexed models instead of uniquely indexed laws, the same coverage convention is applied over those models. Pointwise limits do not establish uniform validity, and asymptotic validity does not guarantee finite-sample precision. Sample sizes, dependence structures and weak-identification sequences are part of each application.

\textbf{Counterfactuals.} $Y_i(a)$ is a potential outcome under a specified intervention, including its timing, implementation and versions. It is not $Y_i$ conditional on voluntarily choosing $a$. A full intervention vector permits interference; shorthand scalar treatment notation does not silently impose no interference. Assignment, selection, attrition, anticipation and measurement must be specified. A claim about an instrument's local effect is not automatically a population policy effect.

\textbf{Economic equilibrium.} A counterfactual model includes best responses, constraints, feasibility, market clearing, state transitions and expectations consistent with the stated information. When several equilibria exist, either a declared selector is an input or the target is set-valued. Comparisons across policy regimes must use the stated selection convention. Reduced-form relationships are not presumed invariant after an equilibrium-changing intervention.

\textbf{Optimization and welfare.} Optimization is over the stated feasible set. If attainment is not established, $\sup$ or $\inf$ and $\varepsilon$-optimal choices replace an asserted maximizer or minimizer. For example, with value $V=\sup_{a\in\mathcal A}W(a)<\infty$, an $\varepsilon$-optimal set is $\{a\in\mathcal A:W(a)\geq V-\varepsilon\}$ for $\varepsilon>0$. Benefits, costs, risk and health or environmental outcomes can be aggregated only using declared units and conversion weights. Interpersonal utility weights, the treatment of fiscal transfers and foreign profits, discounting, and the behavioral welfare criterion are explicit inputs. A comparative-static derivative additionally requires existence and appropriate differentiability of the selected counterfactual.

\textbf{Sources.} Local labels [S1], [S2], and so forth restart in every question. Locators identify the relevant abstract, model, section or result at the level actually checked; a bibliographic or abstract-level verification is not represented as inspection of an entire proof. Working papers are distinguished from later publications and changing versions. Source summaries are paraphrased. All URL checks and editorial formulations refer to the audit date stated above.

\subsection*{Common conventions: Additional-source status and mathematical conventions}

\label{audited:conventions}

Of the 100 supplied ECON entries, 65 distinct problems appear here; 32 close
matches refine earlier entries, and three duplicate questions map to existing
entries. Appendix~\ref{app:audited-crosswalk} lists every destination.
The attachment's P/R/S/U distinctions and source audits are retained. Its
precise mathematical formalizations are editorial unless the individual audit
says otherwise. Candidate subsidy targets ECON-004--006 retain their U flags;
the resolved approval-core question ECON-007 updates OP-0202 above.
The following supplied conventions govern both this part and the attributed
ECON refinements elsewhere in the catalogue, subject to local restrictions.

\textbf{Parameterized questions.} A problem family lists inputs that must be fixed before an instance is evaluated: population, horizon, policies, information, data law, model class and welfare weights. These are required choices, not quantities the citations are assumed to specify. Undefined applications should not be treated as identified estimands.

\textbf{Indices and integrability.} Sets of agents, firms and regions are finite unless a population measure is explicitly used. \(i,j\) index units, \(r\) regions, \(t\) dates and \(h\) horizons. \(\mathbb E\) and \(\Pr\) refer to the declared population and law; \(\mathbf1\{A\}\) is an indicator. Every displayed expectation and discounted sum needs existence in the stated domain. Infinite horizons use a discount in \((0,1)\) and a convergence condition; finite horizons may use discount one.

\textbf{Optimization and existence.} A displayed \(\max\), \(\min\), or \(\arg\max\) is conditional on attainment. For finite feasible menus this is immediate; general spaces need continuity and compactness/coercivity or another existence argument. Without attainment, the target is the supremum/infimum and arbitrarily near-optimal feasible choices. \(\inf\varnothing=+\infty\). Empty feasibility and an unidentified target are different issues.

\textbf{Potential outcomes.} \(Y_i(z)\) is the outcome under a specified intervention, not the observed conditional mean among people choosing \(z\). In binary comparisons \(\Delta Y_i=Y_i(1)-Y_i(0)\). Specify assignment, treatment versions, compliance, information and observation horizon. Interference is allowed under a full policy vector; compressed exposure notation needs a justified mapping. No interference, consistency and overlap are substantive assumptions, not automatic consequences of notation.

\textbf{Populations and observation.} Fix baseline eligibility and population weights. Conditioning on post-policy employment, survival, relocation or program uptake can change the population being compared. Death is not ordinary loss to follow-up; outcomes truncated by death need a composite convention, joint survival target, or explicitly defined survivor stratum. Fertility-changing interventions need a population functional rather than an unexplained fixed descendant cohort. Measurement and missingness models must be supplied.

\textbf{Identification.} A structural model \(m\in\mathcal M\) implies observable law \(P_m\) and target \(\theta(m)\). Identification means
\[
 P_{m_1}=P_{m_2}\ \Longrightarrow\ \theta(m_1)=\theta(m_2)
 \quad\text{for all }m_1,m_2\in\mathcal M .
\]
Without identification, the target is
\[
 \Theta(P;\mathcal M)=\{\theta(m):m\in\mathcal M,\ P_m=P\}.
\]
Writing \(\theta(P)\) before establishing identification can hide incompatible latent models. Confidence sets additionally need a sampling law, confidence level, asymptotic sequence and uniformity class. Point identification, coverage and informative precision are separate properties.

\textbf{Mediation.} Total effects of randomized assignment are distinct from cross-world natural indirect effects. Belief/effort-channel effects require a meaningful mediator intervention and additional measurement, exclusion or mediation assumptions. Randomization of the initial policy alone does not supply those assumptions. Report bounds or interventional alternatives when the stronger target is unsupported.

\textbf{Equilibrium.} A counterfactual \(W(z)\), \(p(z)\), or \(\mu_t^z\) requires individual optimization, feasibility, government/resource budgets, market clearing and state-transition laws. State equilibrium existence and selection, or report ranges over compatible equilibria. Initial-state integration includes subsequent endogenous transitions; current-state integration must use the policy-dependent distribution. Reduced-form invariance after policy is not assumed.

\textbf{Welfare accounts.} Scores, utility, health, dollars and ecological quantities require explicit conversion before addition. Fees, taxes, subsidies, wages and extortion payments are transfers between parties; distributional weights can make them welfare relevant, but their face value is not automatically a resource loss. State ownership, geographic boundaries, foreign-income weights and fiscal finance. Do not add profits to household welfare already including dividends, or count land capitalization and the underlying benefit twice. A benefit-cost expression must distinguish gross benefits from net welfare and subtract every real cost exactly once.

\textbf{Derivatives and risk.} Log derivatives require positive arguments; local responses require admissible perturbations and specified equilibrium branches. Differentiation under expectations needs regularity/domination, and boundaries may allow only one-sided derivatives. Risk uses a specified law; ambiguity uses a declared nonempty law set on a common history space. Adaptive policies must be nonanticipative. A robust criterion is an editorial choice unless attributed explicitly.

\textbf{Scope overlaps.} Allocation, collective choice, contracts, household bargaining and coordinated policy overlap economics with optimization and game theory. They are retained because they appeared in the original catalogue; no claim of a strictly game-theory-free selection is made.
% END ECONOMICS STATEMENT
\end{document}
